% !PRINTED 320
\weberpage{320}
%<<CONT
\weq{2}{\dfrac{(x-\alpha)(x-\beta)f'(x)}{f(x)} = a(x-\beta) + b(x-\alpha) + (x-\alpha)(x-\beta)\dfrac{f_1'(x)}{f_1(x)}.}
The right-hand side of this expression becomes negative for $x = \alpha$, namely
$a(\alpha - \beta)$, and positive for $x = \beta$, namely $b(\beta - \alpha)$,
and must therefore, while $x$ goes from $\alpha$ to $\beta$, pass through zero
once, or more generally an odd number of times. On the left-hand
side of (2), $(x-\alpha)(x-\beta)$ and $f(x)$ have an unchanged
sign, and hence $f'(x)$ must pass through zero an odd number of
times. From this follows Rolle's theorem:

\gesperrt{XII. If $\alpha$ and $\beta$ are two successive roots
of $f(x) = 0$, then between $\alpha$ and $\beta$ (not counting the
limits) there lie roots of $f'(x) = 0$ in odd number, and hence at
least one.}

Important consequences follow from this for the case in which the equation
$f(x) = 0$ has \gesperrt{only real roots}. If these, ordered by magnitude, are
$\alpha, \beta, \gamma, \dots$, and if $\alpha$ is an $a$-fold, $\beta$ a
$b$-fold, $\gamma$ a $c$-fold root, etc., then $f'(x) = 0$ has at
$\alpha$ an $(a-1)$-fold, at $\beta$ a $(b-1)$-fold, at $\gamma$ a
$(c-1)$-fold root, and between $\alpha$ and $\beta$, between $\beta$ and
$\gamma$, and so on, at least one real root each. But the total number of these
roots is
\[
a + b + c + \cdots - 1 = n - 1,
\]
and since $f'(x)$ is of degree $(n-1)$, this minimal
number is also not exceeded. From this follows the theorem:

% [page alignment mirrored 2026-08-07] The German fragment boundary was corrected: theorem
% XIII's tail, the connecting sentence, and theorem XIV are printed on page 320 and were moved
% back here from i0346. The English follows the corrected German, including the source
% emphasis boundary: the theorem numeral is plain and the statement is letterspaced.
XIII. \gesperrt{If $f(x) = 0$ has only real roots, then
$f'(x) = 0$ also has only real roots, and of these, those that do not
coincide with multiple roots of $f(x)$ are simple roots, which are
separated by the roots of $f(x)$.}

If we apply this theorem repeatedly, we can generalise it as
follows:

XIV. \gesperrt{If $f(x)$ has only real roots, then $f'(x), f''(x),
f'''(x) \dots$ also have only real roots. If one of these functions $f^{(\nu)}(x)$
has an $a$-fold root $\alpha$, and if $a > 1$, then $\alpha$ is an
$(a+\nu)$-fold root of $f(x)$.}
